\chapter{Bayesian rationality in verbal reasoning}
\label{sec:p19-rationality-4-5}

Chater and Oaksford's 2021 chapter treats verbal reasoning as a
knowledge-rich activity of revising beliefs, selecting information and
communicating arguments. Its Bayesian proposal is not merely to replace
true and false with numbers. It changes what information is relevant,
what an inference is supposed to accomplish, and how an observed answer
is interpreted. The authors contrast this proposal with the historical
use of logical validity as a psychological benchmark and develop it for
propositional, quantified, inductive, abductive and social reasoning
\cite[pp.~285--294]{ChaterOaksford2021BayesianReasoning}.

Three questions must remain separate: whether a probability assignment
is coherent; whether an updating or search procedure serves a declared
objective under its assumptions; and whether people use that procedure.
A good descriptive fit answers none of the other questions automatically.
The evidence discipline is that of \cref{def:p19-evidence-record}.
The prior, likelihood and observation-law distinctions in
\cref{sec:p15-bayesian-updating-prediction,sec:p15-bayesian-model-comparison}
are inputs to a cognitive account, not a proof of its psychological or
normative adequacy.

\section{A conditional, an instruction and an utterance}
\label{sec:p19-r45-interpretation}

The chapter's initial example is turning a hot tap and obtaining hot
water. With a fixed material conditional and an accepted antecedent,
modus ponens follows; with that same material conditional and a false
consequent, modus tollens follows. These are the fixed-meaning results
in \cref{thm:p19-conditional-benchmark}. Neither result is refuted by
a person treating the hot-tap assertion as a defeasible generalization.

In conversation, a report of not obtaining hot water can suggest that
the speaker tried the tap and the heating failed. That conclusion uses
assumptions about why the speaker made the report. Hearing a report is
not simply learning the event ``no hot water,'' and a normal-operation
generalization is not a material assertion excluding every failure.
The source's more realistic inference therefore changes the interpreted
information rather than producing a countervaluation to classical MT.
The gym-versus-home misunderstanding illustrates another dependency:
speaker and listener must identify the same situation before comparing
their beliefs \cite[pp.~285--287,290]{ChaterOaksford2021BayesianReasoning}.

\begin{hbexample}{def:p19-r45-report-selection}{A stipulated report-selection contrast}
Let \(T\) mean the tap is tried and \(K\) that the heating works.
Independently stipulate \(P(T)=.01\), \(P(K)=.99\), and let
\(W=T\cap K\) be obtaining hot water. Observing \(W^c\) alone gives
\(P(T\mid W^c)=1/9901\). Instead let a report \(M\) occur only
when \(W^c\), with conditional reporting probability 1 if the tap
was tried and \(10^{-6}\) if it was not. Then
\(P(T\mid M)=10000/10099\), and \(P(K^c\mid T\cap W^c)=1\).
The no-hot-water mass is .9901, with tried-tap mass .0001;
the report instead selects masses .0001 and \(.99\times10^{-6}\).
These invented likelihoods demonstrate how a reporting mechanism can
change an inference; they do not estimate how actual speakers behave.
The distinction is not established by labeling a response pragmatic.
\end{hbexample}

The same distinction applies to antecedent strengthening. Material
implication preserves the conclusion after adding a conjunct. A
conditional probability of hot water given a tried tap can fall after
also conditioning on a broken heater. Those are different mathematical
relations. The existing suppression and possibility-premise chapters
retain the actual experiments and interpretation alternatives at
\cref{sec:p19-conditional-reasoning-suppression,sec:p19-conditional-interpretation-experiments}.
The source's first note also reminds readers that philosophical
conditional semantics are not exhausted by material implication.

\section{Coherence is a constraint, not a unique response}
\label{sec:p19-r45-coherence}

For the finite calculations here, propositions denote events under one
declared probability law \(P\); complements and intersections retain
their usual event meanings. Write \(A\) for the antecedent and
\(B\) for the consequent. A conditional-probability interpretation
assigns \(P(B\mid A)\) to the indicative conditional. This is the
source's ``Equation.'' It is an interpretation hypothesis, not a
definition of a new Boolean event whose probability must equal that
ratio under every belief state.

\begin{hbtheorem}{thm:p19-r45-coherence}{proved=annex:p19.r45.coherence}{Sharp uncertain-MP bounds and a different material reading}
Suppose \(P(A)=b\in(0,1)\) and \(P(B\mid A)=a\in[0,1]\).
Then the attainable probabilities of \(B\) are exactly
\[
 [ab,\ ab+1-b].
\]
For every such law,
\[
 1-P(B)\leq (1-a)+(1-b).
\]
This is the uncertainty-sum property for the source's probabilistic-MP
interpretation. By contrast, the material conditional event
\(A^c\cup B\) has probability \(1-b+ab\), equal to \(a\)
exactly when \(a=1\) or \(b=1\), if the latter endpoint is allowed.
If \(b=1\), the consequent probability is \(a\); if \(b=0\),
ordinary positive-event conditioning does not define \(P(B\mid A)\).
\end{hbtheorem}

At the source's \(a=.9,b=.5\), the conclusion range is [.45,.95].
Its uncertainty cannot exceed .55, while the sum of premise
uncertainties is .6. The theorem proves the general property, not just
that example. It does not select a point inside the interval. A person
responding .5 and another responding .9 can both meet this particular
constraint while relying on different knowledge or procedures
\cite[pp.~286--287, note~3]{ChaterOaksford2021BayesianReasoning}.

\begin{hbdefinition}{def:p19-r45-pvalid}{The source's uncertainty-budget criterion}
An argument is p-valid, under a declared interpretation and coherence
framework, when no coherent probability assignment gives its conclusion
an uncertainty greater than the sum of the premises' uncertainties.
An uncertainty here is one minus the corresponding probability.
The quantification is over admissible assignments, not merely over one
fitted assignment or one experimental sample. This formal validity
criterion is not a statistical p-value or a psychological process.
\end{hbdefinition}

The conditional-event accounts discussed by the source can also assign
a conditional assessment at a null antecedent using a richer coherence
framework. That is not the ratio convention in
\cref{def:p15-finite-conditioning}. Applying such an assessment requires
declaring that framework rather than setting \(0/0\) to an arbitrary
number silently. The Equation also does not imply that conditioning
on a further proposition preserves a state-independent conditional
operator. The source's fourth note reports an experimental challenge
to that assumption in an application of a triviality argument; it does
not amount to a proof that every triviality theorem is false. No new
participant count or effect from that uninspected original is asserted
\cite[pp.~286,292, note~4]{ChaterOaksford2021BayesianReasoning}.

The no-life-on-Mars illustration is similarly interpretation-dependent.
In material logic, a false antecedent suffices for an implication.
Under a conditional-probability assessment, knowing the antecedent is
absent does not supply a likelihood for kangaroos on the moon under
that antecedent. Neither the reader's current astronomical knowledge
nor a bare formal implication supplies that missing conditional model.

\subsection{The actual coherence experiments}
\label{sec:p19-r45-coherence-studies}

\begin{hbconvention}{def:p19-r45-study-statistics}{Statistics in the study records}
In the empirical study records, \(F\), \(t\), and \(p\) denote the
reported test statistics and significance values. These are not the
probability parameters or the Jeffrey-update weights defined in the
separate formal models. A nonsignificant contrast does not establish
equivalence, and a significant contrast does not identify a unique
cognitive mechanism.
\end{hbconvention}

Cruz and colleagues' Experiment~1 had 1,140 completed online records
from English-speaking countries, rewarded with EUR~.10. After excluding
duplicate-IP records, invariant responses and implausible reported ages,
871 remained, mean age 35, range 12--78. The reported educational levels
varied and 87\% reported good or very good English. This is a
CrowdFlower sample, not a probability sample of all adult reasoners.
Twenty analyzed respondents were under eighteen; the authors report a
reanalysis excluding them with the same significance pattern.

Participants assessed either separately displayed statements or explicit
premise/conclusion pairs about Linda or a contrasting stereotype,
using percentage confidence. Contents were counterbalanced. Assignment
to booklets was weighted to accommodate their different numbers,
yielding 305 statement and 566 inference participants; it was not an
equal-size two-group design. Experiment~1's six reported forms concern
disjunction introduction and conditional/disjunction transfers, including
negation variants. These forms are not six separate experiments.

The score classified compliance with an interpretation-specific
coherence interval and subtracted the interval's width as a chance
baseline. That baseline assumes uniformly distributed random ratings
over the probability scale. It is not a theorem that human guesses
follow that law, and narrower intervals alter the score even at equal
observed compliance. The reported overall above-chance effect was
\(F(1,869)=885.29,p<.001\). No task effect was significant; that
does not establish equivalence of the tasks or absence of every benefit
from explicit inference \cite[pp.~3--5]{CruzEtAl2015Coherence}.

Experiment~2 tested 48 volunteering University of Orsay students,
native French speakers, mean age 20, range 18--24, mainly biology or
medicine majors. Two equal groups rated statements or explicit
inferences in the university library, in groups of up to four, without
a fixed response deadline. The four reported forms related conjunction
and conditional judgments or eliminated a conjunction, with the Linda
content's prototypicality varied.

The conjunction-to-conditional judgments were coherent above the
authors' baseline, \(F(1,46)=36.83,p<.001\), without a significant
task effect. For conjunction elimination, the content contrast was
\(F(1,46)=33.31,p<.001\); the counter-prototypical condition was
below that baseline, \(F(1,46)=18.37,p<.001\). Explicit inference
did not eliminate the conjunction pattern. Figures show 95\% intervals,
but their endpoints are not printed as exact numerical contrasts.
The published corrigendum swaps the original Figure~2's third and
fourth plotted positions: conclusions here use the source prose and
corrected figure, not the original mislabeled bars
\cite[pp.~5--7]{CruzEtAl2015Coherence}
\cite{CruzEtAl2015CoherenceCorrigendum}.

The conditional-event interpretation makes some responses coherent
that material equivalence would classify differently. This is evidence
relevant to interpretation in these tasks, not permission to choose
whichever semantics maximizes coherence for every new result. Common
language, response scale, selection/exclusion rules and a specified
alternative response process remain part of the comparison.

The 2021 mental-model treatment at \cref{sec:p19-rationality-2-3}
also considers probabilistic judgments and nonvacuous possibilities.
It is not identical to every earlier mental-model formulation used as
a comparator in this literature. A theory name alone does not establish
that all of its versions interpret every conditional materially or
predict the same response. The source/version and the actual task
interpretation must be matched before crediting a theoretical victory.

\section{Updating with certainty, soft evidence and revised links}
\label{sec:p19-r45-updating}

Learning an antecedent for certain under a fixed joint law permits
ordinary conditioning, as in \cref{thm:p15-finite-conditioning-laws}.
Learning something that instead changes its probability can motivate
Jeffrey updating. A new marginal alone, however, does not tell us which
other beliefs to preserve.

\begin{hbdefinition}{def:p19-r45-jeffrey}{A finite binary Jeffrey specification}
Fix a probability law \(P_0\) on a finite event space and
\(0<P_0(A)<1\). Let \(t\in[0,1]\) be a stipulated new
probability of \(A\). Define, for every event \(D\),
\[
 P_1(D)=tP_0(D\mid A)+(1-t)P_0(D\mid A^c).
\]
The specification preserves the old conditional laws within the two
partition cells. This preservation is called rigidity or invariance.
It is an additional assumption about the evidence, not an empirical
property guaranteed by the name Jeffrey.
\end{hbdefinition}

\begin{hbtheorem}{thm:p19-r45-jeffrey}{proved=annex:p19.r45.jeffrey}{The update, its uniqueness boundary and ordinary conditioning}
The displayed formula defines a probability law with \(P_1(A)=t\).
For \(0<t<1\), it preserves both conditional laws and is the unique
law with that marginal and those two conditional laws. At \(t=1\)
it is conditioning \(P_0\) on \(A\); at \(t=0\) it is
conditioning on \(A^c\). The new null cell's conditional law is not
defined by a ratio at either endpoint.

Without rigidity the same new marginal permits every value of
\(P_1(B)\) in [0,1] when no other constraints are imposed on a
four-atom antecedent/consequent space. Consequently the new marginal
alone does not determine the conclusion or a dynamic revision.
\end{hbtheorem}

With \(P_0(B\mid A)=.9\), \(P_0(B\mid A^c)=.1\) and
new \(P_1(A)=.9\), rigidity gives \(P_1(B)=.82\).
The source's hot-tap example instead permits the conditional link itself
to change when failure suggests a broken heater. That possibility is
not a violation of total probability; it changes the inputs to the
Jeffrey formula. Nor does the failure report uniquely determine the
new conditional without a model of the possible causes
\cite[pp.~287--288]{ChaterOaksford2021BayesianReasoning}.

\subsection{Conditional updating experiments}
\label{sec:p19-r45-updating-studies}

Zhao and Osherson's original Experiment~1 recruited seventy Princeton
undergraduates for partial course credit, 45 female, mean age 19.8,
SD 2.1. Forty made probability judgments before and after seeing a
card's color under a dim flashlight; thirty controls judged only after
the light. Group allocation is reported, but a random-assignment
procedure is not specified. The fifty-card deck had twelve blue and
thirty-eight purple cards, with giraffe/hippopotamus images: eight/four
on blue and ten/twenty-eight on purple. Participants saw the complete
deck before the selected card was concealed. The apparently random
draw was deliberately balanced across blue and purple for analysis.

The light exposed color, not the animal. This supplied the authors'
conditional-independence rationale for retaining animal-given-color
probabilities, in contrast with color-given-animal probabilities.
Question order was randomized and respondents could retain or change
an earlier answer. Among the forty experimental respondents, 22 changed
the giraffe-given-blue estimate and eighteen the giraffe-given-purple
estimate. Mean relative movements for the former and its converse were
33\% and 73\%. These denominator-sensitive measures are not percentage
points of accuracy. In the twenty blue-card respondents, mean absolute
total-probability and Jeffrey discrepancies were .09 and .11;
\(t(19)=1.2,p=.24\) did not establish equal errors or an exact
Jeffrey mechanism. The purple group also showed imperfect rigidity.
Twenty additional daylight respondents belong to a separate certainty
comparison, not to the original seventy
\cite[pp.~292--298]{ZhaoOsherson2010Jeffrey}.

Experiment~2 had a new one hundred Princeton undergraduates, 58 female,
mean age 19.5, SD 1.1, for course credit, divided fifty/fifty into
pre/post and post-only conditions. Both assessed a lottery/car-purchase
scenario and an ultimatum-game scenario in counterbalanced order.
Matching the first four lottery numbers was supposed to change the
chance of winning without directly changing car purchase conditional
on a win. Learning a responder had rejected an earlier offer could
instead alter later acceptance even at the same offered amount.
The mean relative movement of acceptance given a high offer was
18.71\%, versus 1.42\% for buying a car given a lottery win,
with paired \(t(49)=3.7,p<.01\). Those are conditional-estimate
movement measures, not probabilities of accepting an offer or making
a purchase. The authors explicitly caution that comparing movement
across domains does not by itself determine whether a change is
normatively warranted.

These are judgments about narrated scenarios, not behavior in an
actual lottery or an incentivized game. The source uses them to test
when rigidity is warranted; it does not support invariance in every
domain. Participant-level records and confidence intervals for the
selected discrepancies are not reconstructed from means
\cite[pp.~299--304]{ZhaoOsherson2010Jeffrey}.

\section{Quantifiers, probabilistic heuristics and prototypes}
\label{sec:p19-r45-quantifiers}

The source represents quantified premises as constraints on a
conditional probability: all gives 1, none 0, some greater than 0,
some-not less than 1. With \(0<\Delta<1/2\), few lies strictly
between 0 and \(\Delta\), and most strictly between
\(1-\Delta\) and 1. These are model interpretations, not an assertion
that every everyday use of few excludes zero or every use of most
excludes all. Conditioning requires a nonnull subject category.

\begin{hbtheorem}{thm:p19-r45-quantifiers}{proved=annex:p19.r45.quantifiers}{The chemist example and the support it requires}
Let \(A,B,C\) denote artist, beekeeper and chemist events.
If \(P(C\cap B)>0\) and \(P(B\cap A^c)=0\), then
\(P(C\cap A)\geq P(C\cap B)>0\). Thus the probability reading
of some chemists being beekeepers and all beekeepers being artists
supports some chemists being artists.

On a finite object domain whose individual atoms all have positive
mass, the probability-one inclusion is an actual set inclusion and
positive intersection is existence. Without that support condition,
an almost-sure inclusion can leave exceptional zero-mass objects.
Nor do two most constraints guarantee the corresponding conclusion:
with \(\Delta=.05\), both conditional probabilities can exceed .95
while \(P(A\mid C)=0\).
\end{hbtheorem}

The chemist--beekeeper--artist chain is a way of displaying premise
links and syllogistic figure. It does not impose conditional
independence or identify a causal mechanism. Treating it as a Markov
chain and multiplying arbitrary transition probabilities would add
assumptions not supplied by the two quantified premises.

The Probability Heuristics Model (PHM) does not ask participants to
solve all of these constraints. Its min heuristic chooses the
conclusion's quantifier using the least informative premise; its max
heuristic controls confidence through the most informative premise.
For the source's all/some example, min favors some and max receives
the all premise. Rarity of predicates motivates its ordering
\[
 \text{all}>\text{most}>\text{few}>\text{some}>\text{none}>\text{some-not}.
\]
The confidence ordering is not identical: some-not and none are
approximately tied. Quantifier choice alone leaves the direction of
the end terms and alternative responses to be specified.

Printed page 288 calls \(I(s)=1/P(s)\) Shannon surprisal. The
inverse probability and Shannon surprisal \(-\log_2 P(s)\) have
the same decreasing ordering in positive \(P(s)\), but are different
numerical scores. Expected scores and fitted confidence rules cannot
be interchanged on that basis. Rarity is itself an environmental
assumption, not a logical consequence of quantifier syntax
\cite[p.~288]{ChaterOaksford2021BayesianReasoning}.

The original PHM paper's generalized-quantifier experiments each
report twenty paid Warwick participant-panel respondents, without prior
experience of psychological reasoning tasks. Experiment~1 used all,
most, few and some-not; Experiment~2 used most, few, some and none.
Each presented 64 two-premise syllogisms in random order, with four
available conclusion quantifiers in a fixed end-term direction.
Respondents could check multiple conclusions or leave all blank, with
no time limit. The response categories are consequently not mutually
exclusive person-level choices. Some probabilistically licensed
conclusions were unavailable in the menus.

Four fitted parameters gave reported aggregate correlations .94 and
.65 and root-mean-square errors .16 and .26 in the two experiments.
These are fits across 256 response-category quantities, not percentages
of participants reasoning correctly or an independent held-out
prediction result. The paper's favorable comparison with a ten-parameter
mental-logic fit involves a different data set. The strong AMFO versus
weaker MFIE fit, conclusion-menu constraints and reported nonsignificant
contrasts remain material limits; no population equivalence follows
\cite[pp.~224--231]{ChaterOaksford1999PHM}.

The Probabilistic Representation Model (PRM) instead constructs a
prototype over eight joint cells of three predicates and their
negations, applies premise constraints, draws a small sample and tests
candidate conclusions in an informativeness order. A sample must be
possible under the premises, but a statement compatible with one
sample need not hold in all premise-compatible structures. The source's
artist/chemist/non-beekeeper sample illustrates this difference
\cite[pp.~288--289]{ChaterOaksford2021BayesianReasoning}.

\begin{hblemma}{thm:p19-r45-sample}{proved=annex:p19.r45.sample}{A finite prototype can miss a counterexample}
If independent draws from a declared prototype law have counterexample
probability \(r\in[0,1]\), the probability that \(n\geq1\)
draws contain no counterexample is \((1-r)^n\).
For \(r=.01,n=7\), it is approximately .9321. Hence a
counterexample-free sample does not establish universal validity when
\(r>0\). This is a stipulated sampling calculation, not an estimated
error rate for the PRM experiments.
\end{hblemma}

Hattori's actual PRM analysis varied sample size from four to twelve
against data from eight earlier experiments. Six minimized RMSD and
seven maximized correlation. These criterion-dependent fitted settings
are not measured participant memory spans. The paper itself discusses
the difference between sampled instances and occupied/chunked cells.
Its generalization to most/few data was inferior to PHM, so the 2021
chapter's short description of a slightly better fit is not a uniform
ranking across every task \cite[pp.~312,315--316]{Hattori2016PRM}.

Hattori's new Experiment~1 recruited 89 Ritsumeikan undergraduates for
course requirements, 55 female and 34 male, mean age 20.2, SD 1.0;
one instruction failure was excluded. Twelve syllogisms were selected
for the predicted inclusion effect, with content intended to distinguish
category prevalence. Booklet order and its reverse were randomly
assigned. Experiment~2 used fifty Kanazawa undergraduates, 35 female
and fifteen male, mean age 20.8, SD 1.0, voluntarily in classwork.
It presented 32 selected computer tasks, but four task types were
excluded because of a recording bug.

The reported AICs for PRM and PHM were 8.080/8.170 and
10.374/10.422; the differences are small and do not identify a unique
representation. The parameterized mental-model comparator required
more parameters in Experiment~1: six versus three for each of the
other two models in that comparison. The fitted coverage parameters
.805 and .779 were compared with earlier-experiment estimates, not
with a newly randomized neutral-content control cohort. No confidence
interval for those fitted settings is supplied in the selected report.
These are content-sensitive model comparisons in particular university
samples, not proof that everyone reasons Bayesianly
\cite[pp.~312--314, Tables~5--6]{Hattori2016PRM}.

\section{Induction, explanation and a precise simplicity bridge}
\label{sec:p19-r45-induction}

Property induction goes beyond deductive entailment. The source
compares attributing an unspecified property from blackbirds to
penguins or worms. Anatomical resemblance may favor penguins, whereas
a poison acquired through a shared environment may favor worms.
Adding ostriches rather than another garden bird can support a wider
property hypothesis. These examples depend on a hypothesis space,
property meaning and a sampling account. Similarity alone does not
prove an inductive conclusion, and diversity need not improve every
inference under every model \cite[p.~289]{ChaterOaksford2021BayesianReasoning}.

The proposed Bayesian treatment of explanatory virtues has a similarly
conditional status. The chapter's single-cause explanation of fatigue
and weight loss is an illustration of integration, not a diagnosis or
an estimate of medical prevalence. The existing original Lombrozo
task at \cref{sec:p19-explanatory-preferences} measures the number of
supplied causes and distinguishes satisfying from likely. That measure
does not become coding length by relabeling it. The source attributes
broader integration and diversity findings to other studies without
supplying a common recruited sample or effect uncertainty; no new
numerical record of those uninspected originals is asserted here.

\begin{hbtheorem}{thm:p19-r45-description}{proved=annex:p19.r45.description}{What a description-length prior actually implies}
On a nonempty finite hypothesis set, assign each hypothesis \(h\)
a finite code length \(L(h)\geq0\) and prior
\(\pi(h)=2^{-L(h)}/Z\), where \(Z=\sum_g2^{-L(g)}>0\).
Let evidence \(D\) have likelihood \(\ell(h)\in[0,1]\),
with positive prior predictive probability. A maximum-posterior
hypothesis minimizes
\[
 L(h)-\log_2\ell(h),
\]
with zero likelihood assigned infinite objective. Minimizing description
length alone is equivalent to this decision only when the candidate
likelihoods are the same positive value, or when an additional argument
shows the likelihood differences cannot change the ordering.
\end{hbtheorem}

For code lengths one and three, and likelihoods .01 and 1, the
longer explanation has posterior probability \(25/26\).
Simplicity can motivate a prior without overriding evidence. The
coding language, hypothesis set and observation likelihood are declared
inputs. The chapter's proposed connection between explanatory virtue
and Bayesian reasoning therefore requires these bridges; empirical
preferences cannot establish them by themselves
\cite[pp.~289--290]{ChaterOaksford2021BayesianReasoning}.

\section{Information search is not complete verification}
\label{sec:p19-r45-search}

The source's blueberry/blue-hair example asks how to investigate a
possible relationship efficiently. Looking among people with rare blue
hair can be diagnostic of a hypothesis even though it is not the
complete counterexample search for a material universal rule. The
contrast is between different inquiry objectives, not an assertion that
confirming observations deductively verify a universal generalization.
The four-card result at \cref{thm:p19-r23-cards} retains its explicitly
complete verification goal.

\begin{hbdefinition}{def:p19-r45-information}{A finite information-gain experiment}
Let a hypothesis \(H\) have a finite positive prior \(\pi(h)\).
For an available test \(a\), stipulate a finite outcome set and
likelihoods \(k_a(y\mid h)\) summing to one for each \(h\).
Let \(m_a(y)=\sum_h\pi(h)k_a(y\mid h)\); for outcomes with
\(m_a(y)>0\), Bayes gives posterior \(\pi_a(h\mid y)\).
Use base-two entropy \(S(\pi)=-\sum_h\pi(h)\log_2\pi(h)\),
with \(0\log_2 0=0\). Expected information gain is
\[
 G(a)=S(\pi)-\sum_{y:m_a(y)>0}m_a(y)S(\pi_a(\cdot\mid y)).
\]
The test's costs and the utility of downstream decisions are not
automatically included in this entropy objective.
\end{hbdefinition}

\begin{hbtheorem}{thm:p19-r45-information}{proved=annex:p19.r45.information}{The expectation, information neutrality and its boundary}
For the declared finite experiment,
\[
 G(a)=\sum_{h,y:\pi(h)k_a(y\mid h)>0}
 \pi(h)k_a(y\mid h)\log_2
 \frac{k_a(y\mid h)}{m_a(y)}\geq0.
\]
It is zero exactly when the outcome law is the same for every
positive-prior hypothesis. A positive or confirming outcome with
different likelihoods can change posterior odds, as can a disconfirming
one. Nonnegative expected gain does not require every realized outcome
to reduce entropy, and maximizing gain does not by itself minimize
expected decision loss or acquisition cost.
\end{hbtheorem}

For two equiprobable hypotheses and a binary test with success
probabilities .8 and .2, a success gives posterior .8 for the first,
a failure .2, and expected gain is \(1-S(.8,.2)\), about .2781 bits.
Neither result eliminates the competing hypothesis. A costless test
independent of the hypothesis has zero gain even if an outcome is
verbally describable as a positive case. Priors, hypothesis alternatives,
sampling rules and the objective determine the comparison.

Oaksford and Chater's original 1994 analysis compares a dependence
model with an independence alternative. The visible card face is
conditioned on when computing the distribution of the hidden face:
turning a displayed card is not drawing a completely unobserved card.
The paper notes that displayed faces chosen by an experimenter cannot
automatically be treated as a random sample informing the hypothesis.
Rarity of the antecedent/consequent predicates motivates the search
ranking, but remains an environmental assumption rather than a law
of card syntax \cite[pp.~608--612]{OaksfordChater1994Selection}.

Its empirical work is a rational reanalysis of prior task studies,
not a new population sample recruited for a common experiment. It
uses a response mapping in addition to the information objective:
an offset of .1 is added to card gains and scores are scaled by
their mean to represent competitive selection. Parameter sampling
and constraints on relative antecedent/consequent frequencies support
its comparisons of associations between card choices. These are model
and response assumptions, not consequences of Bayes' theorem alone.
Its prior-study analyses include negated rules, computer-item versions,
therapy and reduced arrays, whose materials and objectives differ.
No newly reconstructed participant count, pooled prevalence, confidence
interval or independent held-out accuracy is supplied from the printed
correlations. The source supports a particular information-search
proposal under these assumptions; it does not establish that the
complete-verification goal is erroneous
\cite[pp.~614,620--621]{OaksfordChater1994Selection}.

\section{Social arguments, slopes and allegedly circular evidence}
\label{sec:p19-r45-argumentation}

Argumentation requires agreement about what a statement means, what
the audience already believes, and whether a source is reliable.
The hot-tap example illustrates inferential work hidden before the
ostensible MP/MT task begins. A Bayesian account can represent those
dependencies, but merely saying ``all background knowledge'' is not
an operational model: the pertinent alternatives, likelihoods and
response policy must still be specified.

The source's slippery-slope examples compare a same-sex-marriage/pet
marriage claim with a voluntary/involuntary-euthanasia claim. They
are source-attributed examples of perceived category distance and
possible consequences, not this handbook's forecasts of present law,
medical practice or public attitudes. A warning requires both a
consequence probability and a valuation; an unwelcome consequence
does not prove a high transition probability. Perceived similarity or
category flexibility may affect judgments without determining actual
policy transitions \cite[pp.~290--291]{ChaterOaksford2021BayesianReasoning}.

\begin{hblemma}{thm:p19-r45-slope}{proved=annex:p19.r45.slope}{The omitted baseline in a simplified warning rule}
Under a declared planning model, let action 1 provide benefit \(b\),
action 0 provide zero, and let an adverse consequence have disutility
\(c\geq0\) with action-specific probabilities \(q_1,q_0\).
Expected-utility difference is \(b-c(q_1-q_0)\).
Action 1 is strictly preferred exactly when that quantity is positive.
The simpler comparison \(b>cq_1\) agrees when \(q_0=0\)
and all other action costs/outcomes are identical. Observational
\(P(Q\mid A)\) need not identify the action-specific planning laws.
\end{hblemma}

For \(b=3,c=10,q_1=.4,q_0=.2\), action 1 gains one unit in
expected utility although the baseline-free rule would reject it.
This constructed contrast exposes an assumption, not a recommended
policy or an estimated slope probability. The review's reported
similarity and acceptability judgments concern particular argument tasks;
they do not identify action-specific causal probabilities or a unique
mental process. Evaluating that empirical comparison requires its actual
sample, allocation, response measures and uncertainty, rather than the
stipulated warning calculation
\cite[pp.~291--292]{ChaterOaksford2021BayesianReasoning}.

The circularity comparison likewise separates a statement's logical
form from evidential novelty. Repeating \(p\) as its own reason
does not provide a new likelihood. A cloud-chamber observation can
instead discriminate theories if its probability differs between
electron and alternative hypotheses under a model of the apparatus.
Using a hypothesis to calculate that likelihood is not the same as
already accepting the hypothesis as an observed premise. The precise
three-centimeter track example is illustrative, not an independent
particle-identification experiment conducted here.

The source compares this scientific case with an appeal to a text's
claimed divine authorship. Its Bayesian contrast turns on alternative
explanations for the report, not a theorem settling theological truth.
If competing hypotheses assign equal likelihood to a report, the
posterior odds do not change by \cref{thm:p15-bayes-event-odds}.
Counting the same report twice as independent evidence would add a
false independence assumption. The review's reports of responses to
apparently circular arguments are distinct from a quantitative appraisal
of their experiments. Task interpretation, sampling, response coding and
uncertainty are needed before comparing mechanisms or generalizing a
response pattern; a Bayesian representation alone supplies none of them
\cite[pp.~291--292]{ChaterOaksford2021BayesianReasoning}.

\section{What the Bayesian proposal establishes and leaves open}
\label{sec:p19-r45-use}

The chapter proposes continuity between deduction, induction,
explanation and argument rather than a rigid separation into a logical
and an associative system. A common probability framework can describe
several task families; it does not prove that a single psychological
mechanism implements all of them. The source's cited signal-detection
dimensionality study is a particular architectural comparison, not a
logical impossibility theorem for dual-process accounts. No new
participant denominator or effect from that uninspected original is
invented here.

\begin{hbmethod}{def:p19-r45-evaluation}{Evaluating a Bayesian account of a reasoning task}
Specify the interpreted propositions, complete information available
to the participant, prior/likelihood or coherence constraints, response
format and intended task goal. Distinguish deriving a sharp range from
predicting a point, conditioning from revising a link, verifying a rule
from choosing informative evidence, and ranking hypotheses from
choosing an action. Identify which rival account would predict a
different response under the same interpreted task. Check predictions
on observations not used to fit or select the account, and retain
sampling, exclusion, measurement and uncertainty limits. A fitted
parameter is not an observed representation without another bridge.
\end{hbmethod}

An actual mathematical proof can require validity with premises and
meanings fixed. An everyday investigation can instead require revision
under uncertain evidence. The source's concluding concession that
validity remains important in special tasks must survive its broader
proposal. Semantics, pragmatics and background knowledge are substantive
parts of the problem, not licenses for an unconstrained model to explain
any answer after the fact. The valid finite illustrations and bounded
empirical records here clarify these distinctions without promoting
fit into proof of rationality
\cite[p.~292]{ChaterOaksford2021BayesianReasoning}.
